\documentclass[10pt,a4paper]{amsart}
\usepackage[margin=19mm]{geometry}
\usepackage{amsmath,amssymb,mathtools}
\usepackage[scr=rsfs,cal=euler]{mathalfa}
\usepackage{xfrac}
\usepackage[unicode,hidelinks,bookmarks=false]{hyperref}
\newcommand{\ie}{i.e.\ }
\newcommand{\cf}{cf.\ }
\newcommand{\resp}{respectively\ }
\newcommand{\Subsection}{\subsection}
\providecommand{\nobreakdash}{\nobreak\hbox{-}}
\setlength{\emergencystretch}{2em}
\multlinegap0pt
\usepackage{amssymb}
\usepackage{float}
\usepackage{tikz}
\newtheorem{theorem}{Theorem}
\title{Eigenvalues of $\boldsymbol{L_\alpha}-$matrices under graph operations}
\author{Gabriel Roberto Silva de Lima, Carla Silva Oliveira,, Jo\~ao Domingos Gomes da Silva Junior}
\date{Source: arXiv:2605.10944v1 (CC BY 4.0)}
\begin{document}
\maketitle

\noindent\emph{Source and licence.} Eigenvalues of $\boldsymbol{L_\alpha}-$matrices under graph operations. Version arXiv:2605.10944v1 (CC BY 4.0), distributed by arXiv under the Creative Commons Attribution 4.0 International licence, http://creativecommons.org/licenses/by/4.0/, as recorded in the arXiv licence metadata for this exact version and shown on the abstract page https://arxiv.org/abs/2605.10944v1. Full licence notice: \texttt{licenses/arxiv-2605.10944-CC-BY-4.0.txt}.\par\smallskip

\noindent\textit{Editorial scope.} Counted unit: the article's theorem theorem (source0.tex lines 765--775). The article's complete original proof of it is delivered with it (source0.tex lines 777--803, one \texttt{proof} environment of 27 lines). No mathematics is written, repaired or re-derived: the statement and the proof are the article's own lines, in the article's own notation, and no retained line is substituted or re-flowed.  Retained uncounted as the source-local context the proof needs: lines 638--640; lines 687--695; lines 757--761.  Nothing was omitted from any retained span, so the package carries no omission record.  No retained line contains a citation, so the wrapper carries no bibliography and no bibliography entry.  No retained line contains a figure environment, an image-inclusion command or drawing code, so no image is delivered and the package declares no asset dependency.  Editorial formatting only, in addition: an AMS \texttt{amsart} preamble that restates the article's own declarations and macros used by the retained lines, namely the environments \texttt{theorem} on separate counters, together with the standard packages those lines need.  The retained environments print in this document as Theorem 1 in order of appearance, whereas the article prints the same items with the numbering of its own full document.  The complete native source of the article as distributed in the arXiv e-print for this exact version is bundled unchanged as \texttt{sources/200330/source0.tex}.
\begin{equation}\label{eq:cartesian}
L_\alpha(G \times H) = L_\alpha(G) \otimes I_{n_2} + I_{n_1} \otimes L_\alpha(H).
\end{equation}

\begin{theorem}\label{teorema35}
Let $H$ be an $r$-regular graph of order $n_2$, and let $G$ be a connected graph of order $n_1$. Then
\[
\bigcup_{i=1}^{n_1} \{ r\,\lambda_i(L_\alpha(G)) \}
\subseteq
\mathrm{Spec}(L_\alpha(G \odot H)).
\]
for $\alpha \in [0,1].$
\end{theorem}

\begin{equation}\label{eq:otimes}
L_\alpha(G \otimes H)
=
L_\alpha(G \times H) + L_\alpha(G\odot H).
\end{equation}

\begin{theorem}
Let $H$ be an $r$-regular graph of order $n_2$, and let $G$ be a connected graph of order $n_1$. For $\alpha \in [0,1]$, we have
\[
\left\{
(r+1)\lambda_i(L_\alpha(G)) + (2\alpha r - r)
\right\}
\subseteq
\mathrm{Spec}(L_\alpha(G \otimes H)),
\]
for $1 \le i \le n_1.$
\end{theorem}

\begin{proof}
Let $x_i$ be an eigenvector of $L_\alpha(G)$ associated with $\lambda_i(L_\alpha(G))$, and let $j_{n_2}$ be the all-ones vector. From~\eqref{eq:cartesian}, we have
\[
L_\alpha(G \times H)(x_i \otimes j_{n_2})
=
\lambda_i(L_\alpha(G)) (x_i \otimes j_{n_2})
+
L_\alpha(H)j_{n_2} \otimes x_i.
\]
Since $H$ is $r$-regular,
\[
L_\alpha(H)j_{n_2} = (2\alpha r - r)j_{n_2}.
\]
Moreover, by Theorem~\ref{teorema35},
\[
L_\alpha(G \odot H)(x_i \otimes j_{n_2})
=
r\,\lambda_i(L_\alpha(G))(x_i \otimes j_{n_2}).
\]
Combining these expressions with~\eqref{eq:otimes}, we obtain
\[
L_\alpha(G \otimes H)(x_i \otimes j_{n_2})
=
\big[(r+1)\lambda_i(L_\alpha(G)) + (2\alpha r - r)\big](x_i \otimes j_{n_2}),
\]
and the resul follows.
\end{proof}

\end{document}
